\section{Comparing the Estimators}
\label{sec:comparison}
The running example shows how all four methods recover the same expected
gradient while using different information from the mixture. Score-function
estimation weights density derivatives by the observed objective. Stratification
samples every component and differentiates their weighted objectives. Implicit
reparameterization moves one sampled point while keeping its standardized value
fixed. Post-stratification uses one point to average the sample derivatives of
all components, retaining a score term for the weights. These differences
determine both the randomness in an update and the work needed to compute it.

For score-function estimation, the objective value scales the entire density
score, making the baseline an important part of the comparison. Stratification
removes the categorical draw by assigning samples to every component; its
remaining randomness comes from sampling within those components. This also
ensures that a component with a small mixture weight receives samples in each
update. Post-stratification instead averages the component derivatives
conditional on one sampled value. Its Rao--Blackwell guarantee compares it with
the corresponding estimator that uses only the selected component's pathwise
and weight-score terms \citep{burroni2026mixture}. Comparing it with full
stratification or implicit reparameterization requires examining the objective,
the mixture, and the sample budget. In the running example, implicit and
post-stratified estimates agree for the means but differ for the mixture weight,
so even within one problem the comparison depends on which parameters are
being optimized.

The available operations also affect the choice.
Table~\ref{tab:estimators} summarizes one update of each method. Score-function
estimation needs sampling and mixture log-density derivatives, making it a
useful baseline for a broad range of component families. Stratification adds
component sampling derivatives and permits separate sample counts for each
component. Post-stratification uses inverse component maps to obtain averaged
sample derivatives with one objective evaluation. Implicit reparameterization
also uses one objective evaluation and carries weight derivatives through the
sample; its additional requirements are conditional CDFs, their derivatives,
and a Jacobian solve.

\begin{table}[htbp]
 \centering
 \caption{Requirements per update. Stratification uses one sample per component;
 the other methods use one mixture sample.}
 \label{tab:estimators}
 \small
 \begin{tabular}{@{}lcl@{}}
  \toprule
  Estimator & Objective evaluations & Required operations \\
  \midrule
  Score-function & $1$ & Mixture log-density derivatives \\
  Stratified & $K$ & Component sampling derivatives \\
  Implicit & $1$ & Conditional CDFs and a Jacobian solve \\
  Post-stratified & $1$ & Inverse component maps and responsibilities \\
  \bottomrule
 \end{tabular}
\end{table}

Objective counts are particularly informative when evaluating the model is
expensive. In that setting, the single-sample methods can spend the saved model
evaluations on additional updates or independent samples. Density computations
also matter: evaluating the ELBO at one point requires the full mixture
density. For diagonal Gaussians, this entails $O(Kd)$ component-density work,
so evaluating the ELBO once in every component entails $O(K^2d)$ such work.
Post-stratification evaluates all component densities and inverse maps at a
single point, while implicit reparameterization adds the standardization
derivatives and solve described above. Vectorization and the relative cost of
the model determine how these operation counts translate into runtime.

Consequently, an empirical comparison should report both gradient variance at a
fixed mixture and optimization progress under matched computational budgets.
Equal objective-evaluation budgets and equal wall-clock budgets answer different
practical questions, especially as the number of components and the event
dimension grow. STL preserves the expected ELBO gradient and changes the
sample-level estimator, so its effect should be evaluated separately for each
method. ReMix will expose the sample budget and STL option through a shared
interface, making these comparisons possible with the same model and optimizer.
